\section{Benchmark Methodology}
\label{sec:benchmarks}

\subsection{Benchmark Suite Overview}

We developed a comprehensive benchmark suite evaluating Miniforge across four dimensions:

\begin{enumerate}
    \item \textbf{Performance:} Throughput, latency, and processing speed
    \item \textbf{Memory:} Usage patterns, scaling, and efficiency
    \item \textbf{Context:} Retrieval accuracy and maximum effective context
    \item \textbf{Quality:} Output quality across diverse tasks
\end{enumerate}

\subsection{Performance Benchmarks}

\subsubsection{Token Throughput}

Measures generation speed in tokens per second:

\begin{itemize}
    \item Test prompts: short (100 tokens), medium (500 tokens), long (2000 tokens)
    \item Generation length: 512 tokens
    \item Warmup: 2 iterations
    \item Measurement: 5 iterations, report mean and standard deviation
\end{itemize}

\subsubsection{Time-to-First-Token (TTFT)}

Measures latency until first output token:

\begin{itemize}
    \item Streaming generation to capture first token timing
    \item Report mean, P95, and P99 latencies
    \item Critical for interactive applications
\end{itemize}

\subsubsection{Inter-Token Latency (ITL)}

Measures time between consecutive tokens:

\begin{itemize}
    \item Streaming generation with per-token timing
    \item Report mean and P95 ITL
    \item Target: $<$100ms for responsive feel
\end{itemize}

\subsubsection{Prompt Processing}

Measures speed of processing input context:

\begin{itemize}
    \item Context sizes: 512, 1024, 2048, 4096, 8192, 16384, 32768 tokens
    \item Generate single token to measure processing time
    \item Report tokens per second for each context size
\end{itemize}

\subsection{Memory Benchmarks}

\subsubsection{Memory Overhead by Context}

Measures memory consumption at different context lengths:

\begin{itemize}
    \item Context sizes: 512, 2048, 8192, 32768, 65536, 131072
    \item Measure RSS increase after processing each context
    \item Calculate bytes per token scaling factor
\end{itemize}

\subsubsection{Memory Efficiency Score}

Composite metric combining:

\begin{equation}
\text{Efficiency} = 0.4 \times \text{Theoretical Ratio} + 0.3 \times \text{GC Recovery} + 0.3 \times \text{Context Achievement}
\end{equation}

\subsubsection{Generation Memory Stability}

Monitors memory during extended generation:

\begin{itemize}
    \item 512 token generation
    \item Sample memory every 100ms
    \item Measure peak, mean, and variance
    \item Check for memory leaks
\end{itemize}

\subsection{Context Benchmarks}

\subsubsection{Needle-in-Haystack}

Tests retrieval of specific information at different context positions:

\begin{itemize}
    \item Context lengths: 1024, 4096, 16384, 65536, 131072
    \item Depths: 0\%, 25\%, 50\%, 75\%, 100\%
    \item 8 different needle facts tested
    \item Report overall accuracy and accuracy by position
\end{itemize}

\subsubsection{Multi-Needle Retrieval}

Tests retrieval of multiple facts from same context:

\begin{itemize}
    \item Insert 5 different facts into context
    \item Query for each fact separately
    \item Report individual and aggregate accuracy
\end{itemize}

\subsubsection{Maximum Context Test}

Determines effective maximum context window:

\begin{itemize}
    \item Test progressively larger contexts: 8K, 16K, 32K, 64K, 128K, 192K
    \item Simple retrieval task at each size
    \item First failure indicates maximum
\end{itemize}

\subsection{Quality Benchmarks}

\subsubsection{Question Answering}

Tests factual knowledge retrieval:

\begin{itemize}
    \item 50 questions across domains
    \item Check for correct answer presence
    \item Report accuracy percentage
\end{itemize}

\subsubsection{Mathematical Reasoning}

Tests numerical problem solving:

\begin{itemize}
    \item 30 problems from GSM8K subset
    \item Step-by-step reasoning evaluation
    \item Report solve rate and explanation quality
\end{itemize}

\subsubsection{Summarization}

Tests text compression quality:

\begin{itemize}
    \item 10 documents (news, technical, narrative)
    \item Generate summaries
    \item Evaluate concept coverage and coherence
\end{itemize}

\subsubsection{Instruction Following}

Tests adherence to formatting instructions:

\begin{itemize}
    \item 20 tasks with specific format requirements
    \item JSON output, numbered lists, length constraints
    \item Report instruction following accuracy
\end{itemize}

\subsubsection{Coherence}

Tests consistency in long-form generation:

\begin{itemize}
    \item 300-word story generation
    \item Evaluate logical flow and consistency
    \item Check for contradictions
\end{itemize}

\subsection{Benchmark Infrastructure}

\subsubsection{Automated Test Runner}

The benchmark runner orchestrates test execution and result aggregation:

\textbf{Runner Workflow:}
\begin{enumerate}
    \item Initialize benchmark report with timestamp
    \item Iterate through selected benchmark modules
    \item Execute each runner with the target model
    \item Collect results and store in report
    \item Generate statistical summary
    \item Return complete benchmark report
\end{enumerate}

\textbf{Report Contents:}
\begin{itemize}
    \item Timestamp and system configuration
    \item Individual benchmark results with raw data
    \item Statistical summaries (mean, std dev, percentiles)
    \item Pass/fail indicators against targets
\end{itemize}

\subsubsection{Result Storage}

Results stored as JSON with:

\begin{itemize}
    \item Timestamp and system information
    \item Individual test results with raw data
    \item Statistical summaries
    \item Metadata for reproduction
\end{itemize}

\subsubsection{Visualization}

HTML reports with:

\begin{itemize}
    \item Summary tables
    \item Performance charts
    \item Comparison views
    \item Raw data export
\end{itemize}

\subsection{Statistical Methodology}

\subsubsection{Sample Size}

\begin{itemize}
    \item Minimum 5 iterations per measurement
       \item Discard first iteration as warmup
    \item Report mean, standard deviation, and confidence intervals
\end{itemize}

\subsubsection{Outlier Handling}

\begin{itemize}
    \item Identify outliers using IQR method
    \item Investigate and report outliers
    \item Optionally exclude if due to external factors
\end{itemize}

\subsubsection{Reproducibility}

\begin{itemize}
    \item Fixed random seeds where applicable
    \item Document system configuration
    \item Container-based testing for consistency
\end{itemize}
